\problemname{Song Lyrics}
It is well known that the information content of modern song lyrics is not particularly
high.\footnote{http://en.wikipedia.org/wiki/The\_Complexity\_of\_Songs}
We can represent a text using a collection of variables, where each
variable either corresponds to a character string or a concatenation of two
previous variables. The final text is then given by the value of the last
variable.

The competition management now wants to know, for $Q$ different values of $R$, what the $R$th character in
the song lyrics is.

\section*{Input}

The first line contains two integers $N$ ($1 \leq N \leq 500\,000$) and $Q$ ($1 \leq Q
\leq 80\,000$).

Then follow $N$ lines, each containing one of the following two alternatives:

\begin{itemize}
\item A zero followed by a word: \texttt{0 <a word>} (at most $10$ characters in
      the word, only \texttt{a-z}) if the variable represents a simple word.
\item Two integers A and B, the numbers of the concatenated strings
      ($1 \leq A, B < $ current line number). This is thus a word created
      from two merged previous words.
\end{itemize}

Then follow $Q$ lines with one integer $R$ per line ($1 \leq R \leq \min(10^{18},$ length of
the string$))$, the positions of the characters we are interested in.

\section*{Output}

Output the $Q$ requested characters on a single line.

\section*{Scoring}
Your solution will be tested on a set of test groups.
To earn points for a group, you must pass all test cases in that group.

\noindent
\begin{tabular}{| l | l | l |}
  \hline
  \textbf{Group} & \textbf{Points} & \textbf{Constraints} \\ \hline
  $1$   & $20$       & $N \leq 5\,000$, $Q \leq 10\,000$, the length of the strings formed is less than $10^{18}$ characters \\ \hline
  $2$   & $60$       & The length of the strings formed is less than $10^{18}$ characters \\ \hline
  $3$   & $20$       & No additional constraints. \\ \hline
\end{tabular}

\section*{Explanation of Sample 1}
We first get the word ``\texttt{hej}''. Then comes a line that merges
the word with itself, so we now have ``\texttt{hejhej}''. Characters $3$ and $4$ in the string are ``\texttt{jh}''.
